\documentclass[11pt]{article}
\usepackage[utf8]{inputenc}
\usepackage[T1]{fontenc}
\usepackage{lmodern}
\usepackage[brazilian]{babel}
\usepackage[margin=1in]{geometry}
\usepackage{titlesec}
\usepackage{enumitem}
\usepackage{hyperref}
\usepackage{parskip}

\hypersetup{
  colorlinks=true,
  urlcolor=black,
  linkcolor=black,
}

\pagestyle{plain}
\setlength{\parindent}{0pt}
\setlength{\parskip}{0.35em}

\titleformat{\section}
  {\normalsize\bfseries\MakeUppercase}
  {}{0em}{}[\vspace{-0.35em}\hrule\vspace{0.85em}]

\titlespacing*{\section}{0pt}{1.6em}{1.0em}

\newcommand{\cvsubsection}[1]{\vspace{0.9em}\textit{#1}\par\vspace{0.35em}}

\newcommand{\entry}[2]{\noindent\textbf{#1}\hfill #2\par\vspace{0.15em}}
\newcommand{\entryplain}[2]{\noindent#1\hfill #2\par\vspace{0.15em}}
\newcommand{\detail}[1]{\noindent\hspace*{1.5em}#1\par\vspace{0.15em}}
\newcommand{\detailem}[1]{\noindent\hspace*{1.5em}\textit{#1}\par\vspace{0.15em}}

\begin{document}

\begin{center}
{\LARGE \textbf{Anderson Henrique da Silva}}\\[0.5em]
Centro de Estudos da Metrópole (CEM)\\
Universidade de São Paulo --- USP Leste (EACH)\\
São Paulo, SP --- Brasil\\[0.3em]
\texttt{andersonheri@gmail.com} \\[0.15em]
\href{https://www.ahenriquecp.com}{www.ahenriquecp.com} \quad|\quad
\href{http://lattes.cnpq.br/7628365040600511}{Lattes} \quad|\quad
\href{https://orcid.org/0000-0002-1842-2725}{ORCID} \quad|\quad
\href{https://osf.io/ysj2q/}{OSF} \quad|\quad
\href{https://github.com/andersonheri}{GitHub}
\end{center}

\vspace{1em}

\section{Vínculos}

\entry{Pesquisador de Pós-Doutorado}{2025--atual}
\detail{Centro de Estudos da Metrópole (CEM), Universidade de São Paulo. Bolsa FAPESP, proc.\ 2025/15250-0.}

\entry{Pesquisador-Convidado}{2025--atual}
\detail{Instituto de Pesquisa Econômica Aplicada (IPEA), Brasília. Projeto: \textit{Transferências Intergovernamentais no âmbito do Federalismo Brasileiro} (PNPD).}

\entry{Assistente de Pesquisa}{2024--2025}
\detail{Instituto de Pesquisa Econômica Aplicada (IPEA), Brasília.}

\entry{Professor Colaborador}{2023--2025}
\detail{Mestrado Profissional de Sociologia em Rede Nacional (ProfSocio), Universidade Federal do Vale do São Francisco (UNIVASF).}

\entry{Professor Substituto de Ciência Política}{2022--2024}
\detail{Universidade Federal do Vale do São Francisco (UNIVASF), Juazeiro/BA.}

\section{Formação}
\cvsubsection{Titulação}

\entry{Doutorado em Ciência Política}{2025}
\detail{Universidade Federal de Pernambuco (UFPE), Brasil.}
\detailem{Tese: ``Missão: Imprevisto --- Governança, Capacidades Estatais e Desempenho Municipal na Gestão de Risco e Desastres no Brasil.''}
\detail{Orientação: Dalson Britto Figueiredo Filho. Áreas: Estado e Governo; Estudos do Poder Local; Relações Intergovernamentais.}

\entry{Mestrado em Ciência Política}{2021}
\detail{Universidade Federal de Pernambuco (UFPE), Brasil.}
\detailem{Dissertação: ``Desastres naturais, recursos federais de emergência e eleições municipais no Brasil (2012).''}
\detail{Orientação: Gabriela da Silva Tarouco. Coorientação: Davi Cordeiro Moreira.}

\entry{Graduação em Ciência Política}{2018}
\detail{Universidade Federal de Pernambuco (UFPE), Brasil. Orientação: Mariana Batista da Silva.}

\entry{Sanduíche, Relações Internacionais}{2016}
\detail{Universidade de Coimbra, Portugal. Orientação: Carmen Amado Mendes.}

\cvsubsection{Formação complementar}

\entryplain{\textbf{R2}, Research Transparency and Reproducibility Training (32h) --- University of California, Berkeley.}{2024}

\entryplain{\textbf{Survey Data Analysis} (200h) --- University of Essex, Reino Unido.}{2019}

\entryplain{\textbf{Oficina de Análise Legislativa} (6h) --- Universidade Federal de Minas Gerais (UFMG).}{2022}

\section{Bolsas e Auxílios}

\entryplain{Bolsa de Pós-Doutorado FAPESP (proc.\ 2025/15250-0)}{2025--atual}
\entryplain{Bolsa de Doutorado FACEPE}{2021--2025}
\entryplain{Bolsa de Mestrado FACEPE}{2019--2021}
\entryplain{Bolsa de Iniciação Científica CNPq (PIBIC)}{2014--2017}
\entryplain{Bolsa Santander Universidades --- Ibero-Americanas}{2015}
\entryplain{Bolsa PET/MEC --- Programa de Educação Tutorial, UFPE}{2014--2015}

\section{Prêmios e distinções}

\entryplain{Tese de doutorado indicada para premiação pela banca, PPGCP/UFPE}{2025}
\entryplain{Research Transparency and Reproducibility Training (RT2), UC Berkeley}{2024}
\entryplain{Bolsa TSE Summer School, Toulouse School of Economics}{2023}
\entryplain{Menção Honrosa no GT de Instituições Políticas (1º lugar), VII FBCP}{2022}
\entryplain{3º melhor artigo geral, VII Fórum Brasileiro de Pós-Graduação em Ciência Política, UFMG}{2022}
\entryplain{Bolsa Global South Academic Network}{2019}
\entryplain{Voto de Aplausos, Câmara Municipal do Recife}{2017, 2019}

\section{Publicações}
\cvsubsection{Artigos aceitos (no prelo)}
\begin{enumerate}[leftmargin=0.3in,itemsep=0.4em,parsep=0pt]
\item Batista, M.; \textbf{Henrique, A.}; Porto, L. ``Investimento ou aprendizado? Incentivos da reeleição sobre a promoção de capacidades estatais nos municípios brasileiros.'' \textit{Revista Brasileira de Ciências Sociais} (no prelo, 2026).
\item Palotti, P.\ L.\ M.; Linhares, P.\ T.\ F.; \textbf{Henrique, A.} ``Capacidades estatais municipais e o desempenho na gestão de transferências intergovernamentais da União.'' \textit{Texto para Discussão (IPEA)} (no prelo, 2026).
\end{enumerate}

\cvsubsection{Artigos publicados em periódicos}
\begin{enumerate}[leftmargin=0.3in,itemsep=0.4em,parsep=0pt]
\item \textbf{Henrique, A.}; Batista, M. ``Politização dos desastres naturais: alinhamento partidário, declarações de emergência e a alocação de recursos federais para os municípios no Brasil.'' \textit{Opinião Pública}, 26, 522--555, 2020.
\item Silva, D.\ P.\ A.; Figueiredo Filho, D.\ B.; \textbf{Henrique, A.} ``O poderoso NVivo: uma introdução a partir da análise de conteúdo.'' \textit{Política Hoje}, 24, 119--134, 2015.
\item Figueiredo Filho, D.\ B.; Santos Filho, R.\ P.; Rocha, E.\ C.; Silva Jr., J.\ A.; Silva, D.\ P.\ A.; \textbf{Henrique, A.}; Silva, L.\ E.\ O. ```Foi de morte matada': homicídios no Brasil em perspectiva comparada.'' \textit{Sistema Penal \& Violência}, 7, 6--17, 2015.
\item Figueiredo Filho, D.\ B.; Rocha, E.\ C.; Paranhos, R.; Silva Jr., J.\ A.; \textbf{Henrique, A.}; Silva, D.\ P.\ A.; Silva, L.\ E.\ O. ``Análise fatorial garantida ou o seu dinheiro de volta: uma introdução à redução de dados.'' \textit{Revista Eletrônica de Ciência Política --- RECP}, 5, 185--211, 2014.
\end{enumerate}

\cvsubsection{Capítulos de livros}
\begin{enumerate}[leftmargin=0.3in,itemsep=0.4em,parsep=0pt]
\item \textbf{Henrique, A.}; Batista, M.; Porto, L.; Palotti, P.\ L.\ M. ``O Estado em números: construindo bons índices e interpretações de capacidade estatal.'' In: Gomide, A.; Marenco, A.; Oliveira, V.\ E.\ (Orgs.), \textit{Reformar o Estado: o que fazer?}. Porto Alegre: Jacarta/ENAP, 2026, p.\ 224--238.
\item \textbf{Henrique, A.}; Soares, L.\ A.\ C. ``Governança do imprevisto: indicadores, diagnósticos e janelas de oportunidade na gestão de risco e desastres.'' In: Gomide, A.; Marenco, A.; Oliveira, V.\ E.\ (Orgs.), \textit{Reformar o Estado: o que fazer?}. Porto Alegre: Jacarta/ENAP, 2026, p.\ 312--329.
\item \textbf{Henrique, A.}; Braga, D.\ D.\ A. ``Programa Brasil Alfabetizado: uma avaliação da teoria do programa.'' In: Nascimento, P.; Barros, A.\ T.\ D.\ L.\ (Orgs.), \textit{Ciência Política: uma proposta educativa --- Políticas públicas, v.\ 2}. Campina Grande: Eduepb, 2023, p.\ 139--159.
\item Silva, A.\ H.; Batista, M. ``A politização dos desastres naturais: alinhamento partidário e a alocação de recursos federais de emergência para os municípios no Brasil.'' In: \textit{IV Simpósio de Direito Eleitoral do Nordeste}. Campina Grande: Abradep, 2022, p.\ 417--446.
\end{enumerate}

\cvsubsection{Textos em jornais e revistas}
\begin{enumerate}[leftmargin=0.3in,itemsep=0.4em,parsep=0pt]
\item Barros, A.\ T.\ D.\ L.; \textbf{Henrique, A.} ``Preocupémonos por la democracia en Brasil.'' \textit{El País}, 8 jun.\ 2022.
\item \textbf{Henrique, A.} ``O que há de político nos desastres naturais no Brasil?'' \textit{Estadão}, 19 fev.\ 2022.
\end{enumerate}

\section{Palestras e apresentações}

\entryplain{\textit{A política de gestão de risco e desastres no Brasil: quem, quando e como?} --- VI ENEPCP.}{2025}
\entryplain{\textit{Transparência e replicabilidade científica na escrita acadêmica.}}{2024}
\entryplain{\textit{Política e política pública da epidemia de Covid-19 no Brasil} (com Perich, R.; Figueiredo Filho, D.).}{2023}
\entryplain{\textit{Desastres naturais, recursos federais de emergência e eleições municipais} --- VII FBCP.}{2022}
\entryplain{\textit{Eleições no Brasil frente à crise democrática.}}{2022}
\entryplain{\textit{Teste de hipóteses e técnicas de comparação de médias aplicados à saúde.}}{2019}
\entryplain{\textit{O efeito do alinhamento político nas transferências federais de emergência no Brasil.}}{2019}
\entryplain{\textit{The map of public-private partnerships in sanitation in Brazil} --- 2nd LAPolMeth Meeting.}{2018}
\entryplain{\textit{More resources or a better local implementation?} --- VIII ALACIP.}{2015}

\section{Docência}

\cvsubsection{Pós-graduação --- Mestrado Profissional de Sociologia (ProfSocio, UNIVASF)}
\entryplain{Metodologia de Pesquisa}{2024, 2025}
\entryplain{MEDIALIT: Letramento para as Mídias Digitais e Ensino de Sociologia na Escola}{2025}

\vspace{0.9em}
\cvsubsection{Graduação --- Ciências Sociais e Engenharias (UNIVASF)}
\entryplain{Métodos Quantitativos e Interpretação Estatística}{2023}
\entryplain{Métodos Qualiquantitativos}{2023}
\entryplain{Políticas Públicas}{2022, 2023}
\entryplain{Teoria Política I e III}{2022, 2023}
\entryplain{Federalismo, Descentralização e Poder Local}{2024}
\entryplain{Introdução aos Trabalhos Científicos}{2024}
\entryplain{Sociologia (para Engenharia Civil, de Produção e Mecânica)}{2023, 2024}
\entryplain{Bases Sociológicas da Psicologia}{2024}
\entryplain{Bases Filosóficas da Psicologia; Introdução à Filosofia}{2024}
\entryplain{Fundamentos de Antropologia e Sociologia}{2022}

\vspace{0.9em}
\cvsubsection{Estágio docência --- UFPE}
\entryplain{Políticas Públicas I (com M.\ Batista)}{2022}
\entryplain{Seminário de Pesquisa (com A.\ Steiner)}{2020}
\entryplain{Instituições Políticas I (com M.\ Batista)}{2019}
\entryplain{Antropologia (com P.\ Schörder)}{2018}
\entryplain{Métodos Quantitativos I; Teoria Política I (com D.\ B.\ Figueiredo Filho)}{2015--2016}

\section{Grupos e redes de pesquisa}

\entryplain{Centro de Estudos da Metrópole (CEM), USP --- Pós-Doutorando}{2025--atual}
\entryplain{INCT sobre Qualidade de Governo e Políticas para o Desenvolvimento Sustentável (INCT-QualiGov)}{2020--atual}
\entryplain{Grupo Instituições, Governo e Políticas Públicas, DCP/UFPE}{2013--atual}
\entryplain{Polikt --- Instituições, Participação e Cultura Política, UNIVASF}{2022--atual}
\entryplain{Consórcio Crise e Política Local (CPL), UFABC/UFPR/UFPE}{2023--atual}
\entryplain{DERN/GSAN Fellows' Network, University of Essex (Membership Team)}{2021--2022}

\section{Atividade profissional}
\cvsubsection{Corpo editorial}
\entryplain{\textit{Policy Making --- Latin American Public Policy Review}}{2025--atual}

\vspace{0.9em}
\cvsubsection{Parecerista ad hoc}
\textit{Opinião Pública}; \textit{Política \& Trabalho} (UFPB); \textit{Agenda Política}; \textit{Journal of Public Policy and Government}; \textit{Política Hoje} (UFPE); \textit{E-Legis --- Revista Eletrônica do PPG da Câmara dos Deputados}; Prêmio ANPOCS de Teses e Dissertações (2026); Prêmio SEBRAE Prefeitura Empreendedora (2024).

\section{Concursos públicos (aprovações)}

\entryplain{UNIRIO --- Ciência Política / Desenho de Pesquisa. 3º lugar geral; 1º lugar nas vagas PPPIQ.}{2026}
\entryplain{UFAL --- Ciência Política. 3º lugar.}{2025}
\entryplain{UERJ --- Política Internacional / Métodos de Pesquisa em RI. 2º lugar.}{2024}
\entryplain{UFCG --- Ciência Política. 5º lugar.}{2024}
\entryplain{UNIVASF --- Ciência Política (Professor Substituto). 1º lugar.}{2022}

\section{Competências e idiomas}
\cvsubsection{Software}
R, Stata, SPSS, NVivo, \LaTeX, QGIS.

\cvsubsection{Idiomas}
Português (nativo); inglês (fluente); espanhol (intermediário).

\vspace{1.2em}
\begin{center}
\footnotesize Última atualização: \today
\end{center}

\end{document}
